\section{Introduction}\label{sec:intro}

A standard way to certify that a trained network generalizes is to hold out labelled data. Without held-out data, the classical tools (VC dimension, parameter counting, norm-based bounds) are vacuous for networks with more parameters than training examples. They bound every network in a huge class, not the one that was trained.

This paper studies a reduction that targets the trained network itself. Suppose that some network $g$ with a \emph{short description} (a code of $L$ bits fixed before seeing the data) agrees with the trained network $f$ on most inputs. Then
\[
\err(f)\ \le\ \err(g)+\Pr_x[f(x)\ne g(x)],
\]
and since $g$ ranges over at most $2^L$ descriptions, a Chernoff--Occam union bound over all of them bounds $\err(g)$ by its training error plus a term that is about $L\ln2/m$ when $g$ has no training errors (and of order $\sqrt{L/m}$ otherwise). Bounding generalization reduces to compressing the network, provided one can certify the closeness term $\Pr[f\ne g]$. That term concerns the input distribution, which is exactly what the training data does not reveal without depending on $f$.

We consider three regimes, in decreasing order of what is known about the input distribution, and ask three questions. Has this been done? Can a compression scheme give interesting bounds on real networks, beating parameter counting? Does the random initialization have to be described, which would make the description long?

\paragraph{Answers in brief} (\cref{tab:headline}).
\begin{itemize}
\item \emph{Prior work.} Compression bounds for the compressed network $g$ are well developed \cite{arora2018stronger,zhou2019nonvacuous,lotfi2022pac}. The reduction to the original $f$ with the closeness term estimated on \emph{unlabeled} data goes back to \citet{kaariainen2005generalization} and was recently carried out for modern networks by \citet{bazinet2026bound}, with non-vacuous bounds on MNIST and CIFAR-10. With no extra data, the deterministic derandomization bounds of \citet{neyshabur2018pacbayes} and \citet{nagarajan2019deterministic} bound the original network but are vacuous. The exception we know of is \citet{biggs2022nonvacuous}, who obtain non-vacuous training-data-only bounds for one-hidden-layer networks with erf or GELU activations, via an exact identity between such a network and a stochastic majority vote.
\item \emph{Known input distribution} (\cref{sec:exp-bool,sec:exp-anyD}). On Boolean tasks with uniform inputs, $\Pr[f\ne g]$ is free. A 289{,}281-parameter MLP trained on $n=50$ bits compresses into a 28--100-bit truth table, linear threshold function or small quantized network. The resulting bound on $f$ is within $0.005$ of its test error at $m=16000$. At smaller $m$ the Occam term makes it looser, but it stays informative: $0.30$ against $0.29$ for a partially trained parity network, $0.34$ against $0.17$ for a partially trained 6-junta, and $0.17$ against $0.14$ for a network that interpolates 10\% label noise. Parameter-counting and VC bounds are 1.0 throughout. For $n=20$, once $f$ has learned the target, the compressed networks agree with $f$ on all but at most 8 of the $2^{20}$ inputs. Patching those exceptions gives bounds that hold for \emph{every} input distribution (0.0065 for \texttt{parity3} at $m=4000$).
\item \emph{Training data only} (\cref{sec:transfer,sec:exp-t3}). We prove a \emph{dithered transfer theorem}. Compress $f$ by subtractive dithered quantization with public randomness, so that $g-f$ is independent uniform noise in weight space; compute from $g$ alone a per-input certificate that $f$ and $g$ classify alike; and transfer the fraction of certified training points to the population with the same Occam bound. The result is a deterministic bound on the original $f$ (the stored weights, evaluated in float64) with no validation or unlabeled data. We prove the per-input certificate for one-hidden-layer ReLU networks. The bound is non-vacuous for narrow one-hidden-layer MNIST networks: 0.44 for 8 hidden units, where the VC bound with $d$ equal to the number of parameters gives 0.71. It is also non-vacuous for standard-width networks whose first-layer update is sparse in a DCT basis: 0.31 for 100 hidden units with 1\% nonzero update coefficients, where coding the network exactly gives 0.62 (\cref{tab:t3sparse}). It becomes vacuous for wide dense networks. The certificate needs $f$'s weights described to the precision its margins require, and that costs bits in proportion to the width. Worst-case versions of the same idea fail by a factor $\sqrt d$ (\cref{thm:certified}).
\item \emph{Initialization} (\cref{sec:init-theory,sec:exp-init}). The initialization costs nothing: any data-independent randomness may be used by the code (\cref{lem:init}). Function-space codecs never refer to it, and weight-space codecs code the update relative to it. We measure how coding relative to the initialization or to zero, the initialization scale and a deterministic initialization affect the description length. We give an init-kernel codec (\cref{prop:kernel}). On lazily trained Boolean networks with majority targets its best rank is $1+n$, the dimension of the degree-$\le1$ Walsh space; this suggests, as a heuristic we have not tested across $n$, that its cost grows like the $n^k$ sample complexity of kernel methods for degree-$k$ targets. In that picture, feature learning is what makes trained networks short to describe.
\end{itemize}

\begin{table}[t]
\centering\small
\begin{tabular}{l l l l r r r r}
\toprule
task & network $f$ & extra data & data & test err.\ $f$ & bound on $f$ & bits & 32-bit Occam \\
\midrule
\texttt{parity3} & MLP, 289,281 params & known $D_X$ & 50-bit, $m=4000$ & 0.0000 & 0.007 & 35 & 1.00 \\
\texttt{parity3}, partly learned & MLP, 289,281 params & known $D_X$ & 50-bit, $m=2000$ & 0.2876 & 0.304 & 35 & 1.00 \\
\texttt{junta6} & MLP, 289,281 params & known $D_X$ & 50-bit, $m=4000$ & 0.0223 & 0.042 & 102 & 1.00 \\
\texttt{maj5}, 10\% label noise & MLP, 289,281 params & known $D_X$ & 50-bit, $m=16000$ & 0.1390 & 0.174 & 53 & 1.00 \\
\texttt{parity3} & MLP, 273,921 params & none (any $D$) & 20-bit, $m=4000$ & 0.0000 & 0.006 & 33 & 1.00 \\
MNIST & 784-8-10 & none & $m=60000$ & 0.0697 & 0.439 & 23,490 & 0.94 \\
MNIST & 784-16-10 & none & $m=60000$ & 0.0512 & 0.553 & 43,658 & 0.99 \\
MNIST & 784-100-10, sparse update (1\%) & none & $m=60000$ & 0.0462 & 0.314 & 17,664 & 1.00 \\
MNIST & 784-256-10, sparse update (0.3\%) & none & $m=60000$ & 0.0441 & 0.399 & 23,244 & 1.00 \\
\bottomrule
\end{tabular}
\caption{Headline results: certified upper bounds on the test error of the trained network itself, at confidence 95\%, against its measured test error and the parameter-counting bound (32 bits per parameter). `Extra data': what the bound uses besides the training set; no row uses labelled held-out data. Chance level is 0.5 for the Boolean tasks and 0.9 for MNIST. Boolean rows average three seeds. `Bits': description length of the compressed or dithered network.}\label{tab:headline}
\end{table}


\paragraph{Scope.} Everything here is a theorem or a reproducible measurement on small networks: MLPs on 20--50-bit Boolean tasks, MNIST and Fashion-MNIST, on a 4-core CPU. Every bound is computed with explicit code lengths, confidence parameters and a charge for the design choices we made on the same data (\cref{sec:experiments}). The code, the result files and an AI referee record are in the repository.
